\documentclass[11pt,a4paper]{article}
\usepackage[utf8]{inputenc}
\usepackage{amsmath,amssymb}
\usepackage{geometry}
\geometry{margin=2cm}
\usepackage{booktabs}
\usepackage{tabularx}
\usepackage[hidelinks,unicode]{hyperref}
\usepackage{algorithmic}
\usepackage{algorithm}

\begin{document}
\begin{center}
	{\large\bfseries Reference Sheet: Metrics, Algorithms \& Definitions\\for Diversity in RAG Retrieval}
\end{center}
\vspace{-1em}
\tableofcontents

\section{Notation}

\begin{tabular}{@{}ll@{}}
	$D$ & Full dataset/candidate pool of $n$ items (chunks)\\
	$S \subseteq D$ & Selected subset, $|S| = k$\\
	$\mathbf{e}_i$ & Embedding vector of item $i$\\
	$u_i$ & Utility (relevance) of item $i$ to the query\\
	$w_{ij} \geq 0$ & Similarity between items $i$ and $j$\\
	$\cos(a,b)$ & Cosine similarity between vectors $a$ and $b$\\
	$\lambda, \alpha, \beta$ & Trade-off parameters between relevance and diversity\\
	$T_{\max}$ & Context budget (max tokens in selected chunks)\\
\end{tabular}

\section{Quick Reference}

\subsection*{Diversity Definitions}

\begin{tabularx}{\textwidth}{@{}l l l X@{}}
	\toprule
	\textbf{Concept} & \textbf{Scope} & \textbf{Approx.\ guarantee} & \textbf{Key property} \\
	\midrule
	Max-min Dispersion & Selected only  & None (NP-hard) & Maximises closest-pair distance in $S$ \\
	Coverage (Facility-Location) & All of $D$ & $(1-1/e)$ greedy & Monotone submodular; respects data distribution \\
	Combined Objective $F$ & All of $D$ & $(1-1/e)$ greedy & Modular utility $+$ submodular coverage \\
	\bottomrule
\end{tabularx}

\subsection*{Selection Algorithms}

\begin{tabularx}{\textwidth}{@{}X l l l l@{}}
	\toprule
	\textbf{Algorithm} & \textbf{Diversity type}   & \textbf{Relevance} & \textbf{Complexity} & \textbf{Guarantee} \\
	\midrule
	Top-$k$ & None (baseline) & Yes & $O(n \log k)$ & — \\
	MMR & Dispersion (local window) & Yes & $O(k n w)$ & — \\
	gMMR & Dispersion (centroid) & Yes & $O(kn)$ & — \\
	FPS & Dispersion (global) & No & $O(kn)$ & — \\
	Greedy Facility-Location & Coverage & Yes & $O(kn)$ & $(1-1/e)$ \\
	Branch-and-Bound & Coverage & Yes & Exp.\ in $k$ & Optimal \\
	\bottomrule
\end{tabularx}

\subsection*{Retrieval Quality Metrics}

\begin{tabularx}{\textwidth}{@{}l l l X@{}}
	\toprule
	\textbf{Metric} & \textbf{Range} & \textbf{Requires}   & \textbf{What it measures} \\
	\midrule
	Hit Rate ($\mathrm{hit}$) & $\{0,1\}$ & Gold aliases & Binary: does context contain the answer? \\
	Facility-Location ($\mathrm{fac}$) & $[0,1]$ & Embeddings & Average representativeness of $S$ over $D$ \\
	Recall@$K$ & $[0,1]$ & Relevance judgments & relevant retrieved / total relevant \\
	Precision@$K$ & $[0,1]$ & Relevance judgments & relevant retrieved / (relevant retrieved + non-relevant retrieved) \\
	nDCG@$K$ & $[0,1]$ & Graded judgments    & Ranking quality with position discount \\
	$F_1$ & $[0,1]$ & Ground-truth answer & Token-level precision/recall of generated answer \\
	DocRec & $[0,1]$ & Gold documents & Fraction of supporting documents selected \\
	FactRec & $[0,1]$ & Gold facts & Fraction of supporting sentences selected \\
	AvgCos & $[0,1]$ & Embeddings & Average query–chunk cosine similarity \\
	\bottomrule
\end{tabularx}

\subsection*{Selection Overlap Metrics}

\begin{tabularx}{\textwidth}{@{}l l X@{}}
	\toprule
	\textbf{Metric}    & \textbf{Range} & \textbf{What it measures} \\
	\midrule
	Jaccard similarity & $[0,1]$ & Set overlap between two selection strategies; $J \approx 1$ means the diversity method changes nothing \\
	GoldConcat & $\geq 0$ & Number of distinct answer aliases present in the concatenated context \\
	Gold/Chunk & $\geq 0$ & Average answer-alias density per selected chunk \\
	\bottomrule
\end{tabularx}

\newpage

\section{Diversity Definitions}

\subsection{Dispersion}

Diversity measured by distances \emph{between selected items only}. Does not consider unselected items.

\textbf{Max-min dispersion:}
\[
	\mathrm{Div}(S) = \min_{i \neq j \in S} \mathrm{dist}(i, j)
\]
Maximizing this is NP-hard (the $k$-dispersion problem). No polynomial-time approximation guarantee for the greedy algorithm.

\textbf{Intuition.} Push selected points as far apart from each other as possible. The ``weakest link'' (closest pair) defines the score. Natural but ignores the structure of unselected data---may place representatives in empty regions of the space while leaving dense clusters uncovered.

\subsection{Coverage (Facility-Location)}

Diversity measured by how well $S$ \emph{represents the entire dataset} $D$.
\[
	g(S) = \sum_{i \in D} m_i(S), \qquad m_i(S) = \max_{j \in S} w_{ij}
\]
\begin{itemize}
	\item $m_i(S)$: how well item $i \in D$ is covered (similarity to its nearest representative in $S$).
	\item $g(S)$: total coverage over all items in $D$.
\end{itemize}
$g$ is \textbf{monotone submodular}: $A \subseteq B \Rightarrow g(A) \leq g(B)$, and for all $A \subseteq B$, $x \notin B$:
\[
	g(A \cup \{x\}) - g(A) \geq g(B \cup \{x\}) - g(B) \qquad \text{(diminishing returns)}
\]

\textbf{Intuition.} Every point in $D$ should have a close representative in $S$. If 90\% of the data clusters in region~A, coverage allocates more representatives there. It respects the data distribution, unlike dispersion which may ignore dense clusters. Key advantage: submodularity gives provable approximation guarantees.

\subsection{Combined Objective}

\[
	F(S) = \lambda \sum_{i \in S} u_i + (1-\lambda)\, g(S)
\]
Sum of a modular function (utility) and a submodular function (coverage). $F$ is monotone submodular. Greedy gives a $(1 - 1/e) \approx 0.632$ approximation.

\textbf{Intuition.} Trade off between picking the most relevant chunks ($\lambda = 1$, standard top-$k$) and picking chunks that collectively represent the whole candidate pool ($\lambda = 0$, pure coverage).

\section{Selection Algorithms}

\subsection{Top-k by Relevance (Baseline)}

Select $k$ items with highest $u_i = \cos(\mathbf{e}_q, \mathbf{e}_i)$, subject to context budget $T_{\max}$. $O(n \log k)$.

\textbf{Intuition.} The simplest strategy: rank by similarity to the query, take the top. Fast but blind to redundancy---if the top chunks all say the same thing, the context is wasted.

\subsection{MMR (Maximal Marginal Relevance)}

Greedy, dispersion-based. At each step, select:
\[
	c^* = \arg\max_{i \in R} \Big[ \alpha \cdot \cos(\mathbf{e}_q, \mathbf{e}_i) - (1-\alpha) \cdot \max_{j \in W} \cos(\mathbf{e}_i, \mathbf{e}_j) \Big]
\]
\begin{itemize}
	\item $R = C \setminus S$: remaining candidates.
	\item $W \subseteq S$: context window (last $w$ selected items).
	\item $\alpha \in [0,1]$: relevance--diversity trade-off.
\end{itemize}
First term rewards relevance; second penalizes similarity to already-selected items. $O(k \cdot n \cdot w)$ per query.

\textbf{Intuition.} ``Pick the next chunk that is relevant but not redundant with what I already have.'' The penalty is against the single most similar already-selected item (pairwise max). The context window $w$ limits computation and makes the penalty local---only the $w$ most recent selections matter.

\subsection{gMMR (Geometric MMR, from DF-RAG)}

Greedy, dispersion-based. Replaces pairwise max with centroid distance, cosine with Euclidean.
\[
	\mathrm{gMMR}(c) = \lambda \cdot \cos(\mathbf{e}_q, \mathbf{e}_c) + (1-\lambda) \cdot \sqrt{2 - 2\cos(\mathbf{e}_c, \mathbf{e}_{c_S})}
\]
where the centroid of the selected set is:
\[
	\mathbf{e}_{c_S} = \frac{1}{|S|} \sum_{s \in S} \mathbf{e}_s
\]
At each step: $c^* = \arg\max_{c \in C \setminus S} \mathrm{gMMR}(c)$.

\textbf{Intuition.} Instead of penalizing similarity to the nearest selected item, gMMR rewards distance from the \emph{center of mass} of all selected items. This is a more global view of diversity: ``go far from where I've already concentrated,'' not just ``avoid my nearest neighbor.'' Euclidean distance on normalized embeddings is $\sqrt{2 - 2\cos}$, grounding diversity in geometric space.

\subsection{FPS (Farthest Point Sampling)}

Greedy, pure dispersion (no relevance component).
\[
	c^* = \arg\max_{i \in R} \Big[ \max_{j \in S} \mathrm{dist}(i, j) \Big]
\]

\textbf{Intuition.} At each step, pick the point that is farthest from all already-selected points. Pure space-filling strategy with no relevance signal. From 3D computer vision (point cloud subsampling). Slower than MMR because it uses Euclidean distance over all selected items rather than cosine over a window.

\subsection{Greedy for Facility-Location Coverage}

Start with $S = \emptyset$. Repeat $k$ times: add $j^* = \arg\max_{j \notin S} \Delta(j \mid S)$, where the marginal gain is:
\[
	\Delta(j \mid S) = \lambda \cdot u_j + (1-\lambda) \sum_{i \in D} \max\big(0,\; w_{ij} - m_i(S)\big)
\]
\begin{itemize}
	\item Second term counts only the \emph{new} coverage $j$ brings.
	\item $m_i(S) = \max_{s \in S} w_{is}$: current coverage of item $i$.
	\item $O(kn)$ per query (update $m_i$ array incrementally).
	\item Provides a $(1-1/e)$ approximation guarantee since $F$ is monotone submodular.
\end{itemize}

\textbf{Intuition.} Each candidate is evaluated by: (1)~its relevance $u_j$, and (2)~the \emph{new} coverage it adds---only counting items in $D$ that it covers \emph{better} than the current best representative. As $S$ grows, the marginal coverage of each new item diminishes (submodularity), naturally preventing redundant selections.

\subsection{Exact Branch-and-Bound for Coverage}

Finds the optimal $S^* = \arg\max_{|S|=k} F(S)$ exactly.

\begin{enumerate}
	\item Nodes are partial sets $S_p$, $|S_p| \leq k$. Priority queue ordered by upper bound.
	\item \textbf{Upper bound} (exploits submodularity):
	      \[
		      \mathrm{UB}(S_p) = F(S_p) + \sum_{j \in R_{(1:t)}} \Delta(j \mid S_p)
	      \]
	      where $R_{(1:t)}$ are the $t = k - |S_p|$ largest marginals among candidates.
	\item \textbf{Lower bound}: any feasible completion (e.g., greedy from $S_p$).
	\item \textbf{Pruning}: if $\mathrm{UB}(S_p) \leq \mathrm{BestValue}$, discard node.
	\item \textbf{Branching}: expand $S_p$ by adding each candidate $e \in D \setminus S_p$, sorted by descending marginal gain.
\end{enumerate}
Worst-case: $O\!\left(\sum_{i=0}^{k} \binom{n}{i} \cdot n^2\right)$. Polynomial in $n$ for fixed $k$, but exponential in $k$. Practical for small $k$ with effective pruning.

\textbf{Intuition.} The greedy solution is near-optimal but not exact. Branch-and-bound explores the combinatorial space of possible selections, using submodularity-derived upper bounds to prune most of the search tree. The key insight: submodularity guarantees that the sum of the $t$ largest single-item marginal gains is an upper bound on the gain from adding \emph{any} $t$ items, because marginal gains can only decrease as the set grows.

\section{Retrieval Quality Metrics}

\subsection{\texorpdfstring{Hit Rate ($\mathrm{hit}_b$, $\mathrm{hit}_c$)}{Hit Rate (hit\_b, hit\_c)}}

\[
	\mathrm{hit}(q) =
	\begin{cases}
		1 & \text{if any gold alias of } q \text{ appears in } \mathrm{concat}(S_q) \\
		0 & \text{otherwise}
	\end{cases}
\]
Averaged over queries: $\mathrm{hit}_b$ (baseline), $\mathrm{hit}_c$ (coverage). $\Delta\mathrm{hit} = \mathrm{hit}_c - \mathrm{hit}_b$.

\textbf{Intuition.} Binary pre-LLM recall: did the selected chunks contain the answer at all? If $\mathrm{hit}=0$, no LLM can produce the correct answer from this context.

\subsection{Facility-Location Coverage Score ($\mathrm{fac}$)}
\[
	\mathrm{fac}(S) = \frac{1}{|D|} \sum_{i \in D} \max_{j \in S} \cos(\mathbf{e}_i, \mathbf{e}_j)
\]
$\Delta\mathrm{fac} = \mathrm{fac}_c - \mathrm{fac}_b$.

\textbf{Intuition.} Average per-item representation quality. High $\mathrm{fac}$ means every candidate chunk has a close representative among the selected chunks. Measures whether coverage-aware selection actually achieves better representativeness.

\subsection{Recall@$K$}

\[
	\mathrm{Recall@}K(q) = \frac{|\{p \in \text{top-}K : p \text{ is relevant}\}|}{|\{p \in \text{corpus} : p \text{ is relevant}\}|}
\]

\textbf{Intuition.} Fraction of all relevant passages captured in the top-$K$. Used with TREC qrels where human relevance judgments are available.

\subsection{nDCG@$K$ (Normalized Discounted Cumulative Gain)}

\[
	\mathrm{DCG@}K = \sum_{i=1}^{K} \frac{2^{\mathrm{rel}_i} - 1}{\log_2(i+1)}, \qquad
	\mathrm{nDCG@}K = \frac{\mathrm{DCG@}K}{\mathrm{IDCG@}K}
\]
$\mathrm{rel}_i$: graded relevance at rank $i$. IDCG: DCG of the ideal ranking.

\textbf{Intuition.} Not just \emph{whether} relevant passages are retrieved, but \emph{how well they are ranked}. Highly relevant passages at early ranks score much more than at later ranks (logarithmic discount). Normalized to $[0,1]$ by the ideal.

\subsection{$F_1$ Score (DF-RAG)}

\[
	P = \frac{|\text{tokens}(\hat{a}) \cap \text{tokens}(a^*)|}{|\text{tokens}(\hat{a})|}, \qquad
	R = \frac{|\text{tokens}(\hat{a}) \cap \text{tokens}(a^*)|}{|\text{tokens}(a^*)|}, \qquad
	F_1 = \frac{2 \cdot P \cdot R}{P + R}
\]
where $\hat{a}$ is the generated answer and $a^*$ is the ground truth; tokens are compared as bags of words.

\textbf{Intuition.} Balances how much of the gold answer the system produced (recall) with how much of the system's output is correct (precision). Used across all DF-RAG benchmarks for end-to-end QA evaluation.

\subsection{DocRec / FactRec (HotpotQA)}

\[
	\mathrm{DocRec}(q) = \frac{|S_q \cap \mathrm{gold\_docs}(q)|}{|\mathrm{gold\_docs}(q)|}
\]

\textbf{Intuition.} Did we select chunks from the \emph{right} documents? In HotpotQA (2 gold docs out of 10), $\mathrm{DocRec}=1.0$ means both supporting documents are represented.

\[
	\mathrm{FactRec}(q) = \frac{|\mathrm{selected\_facts}(q) \cap \mathrm{gold\_facts}(q)|}{|\mathrm{gold\_facts}(q)|}
\]

\textbf{Intuition.} Finer than DocRec: even if the right document is chosen, did we capture the specific \emph{sentences} annotated as supporting facts? This is the most direct measure of whether diversity helps recover complementary evidence.

\subsection{AvgCos (Average Cosine Similarity)}

\[
	\mathrm{AvgCos}(q) = \frac{1}{|S_q|} \sum_{j \in S_q} \cos(\mathbf{e}_q, \mathbf{e}_j)
\]

\textbf{Intuition.} Average relevance of selected chunks to the query. Lower AvgCos means the strategy sacrificed relevance to gain diversity. Useful for diagnosing whether a diversity method is ``paying too much'' in relevance.

\section{Selection Overlap Metrics}

\subsection{Jaccard Similarity}

\[
	J(A, B) = \frac{|A \cap B|}{|A \cup B|}
\]
Applied at chunk level ($\mathrm{Jaccard}_{\text{chunks}}$) or document/domain level ($\mathrm{Jaccard}_{\text{domains}}$).

\textbf{Intuition.} Measures agreement between two selection strategies. $J \approx 1$: the diversity method barely changes the selection (no effect). $J \ll 1$: genuinely different selections. The most diagnostic metric for determining whether a diversity method is actually doing something.

\subsection{Gold-Alias Metrics (TriviaQA)}

In TriviaQA each query $q$ has a set $A(q)$ of \emph{gold aliases}: multiple valid surface forms of the correct answer (e.g., ``teapot dome scandal'', ``teapot dome'', ``tea pot dome affair'', \ldots). The following metrics count how many of these aliases are recovered by the selected context.

\subsubsection{GoldConcat}

For a single query, let $T_q = \mathrm{concat}(S_q)$ be the concatenation of all selected chunks. The per-query score is
\[
	\mathrm{GoldConcat}(q) = \big|\{ a \in A(q) : a \subseteq T_q \}\big|
\]
and the reported metric is the mean over all queries:
\[
	\mathrm{GoldConcat} = \frac{1}{|Q|} \sum_{q \in Q} \mathrm{GoldConcat}(q).
\]

\textbf{Intuition.} A finer-grained version of the binary hit rate: instead of asking ``is \emph{any} alias present?''\ it counts \emph{how many distinct} phrasings of the answer appear in the full selected context. $\mathrm{GoldConcat} = 1.47$ means the context contains on average $1.47$ distinct correct phrasings per query. Hit rate $= 1$ and $\mathrm{GoldConcat} = 1$ are both consistent with barely covering the answer once; higher GoldConcat signals richer, more thorough evidence.

\subsubsection{Gold/Chunk}

For a single query, score each selected chunk $j \in S_q$ individually:
\[
	\mathrm{Gold/Chunk}(q) = \frac{1}{|S_q|} \sum_{j \in S_q} \big|\{ a \in A(q) : a \subseteq \text{chunk}_j \}\big|
\]
and the reported metric is the mean over all queries:
\[
	\mathrm{Gold/Chunk} = \frac{1}{|Q|} \sum_{q \in Q} \mathrm{Gold/Chunk}(q).
\]

\textbf{Intuition.} Measures \emph{per-chunk} answer density rather than aggregate recall. If $\mathrm{Gold/Chunk} \approx 1$, most selected chunks individually mention the answer; if it is near $0$, the answer is present in the concatenated context but concentrated in a single chunk while the rest are answer-free. Comparing GoldConcat and Gold/Chunk together reveals whether a selection strategy spreads answer evidence across chunks or concentrates it.

\section{Key Experimental Parameters}

\begin{tabular}{ll}
	\toprule
	Parameter & Description\\
	\midrule
	$T_{\max}$ & Context budget in tokens. Tighter budgets force harder selection and amplify diversity effects.\\
	$\lambda$ / $\alpha$ & Relevance--diversity trade-off.\\
	$w$ & Context window size (MMR): number of recent selections considered for the diversity penalty.\\
	chunk\_tokens & Chunk size in tokens for document splitting.\\
	stride & Overlap between consecutive chunks (stride $<$ chunk\_tokens implies overlap).\\
	max\_docs & Number of source documents contributing to the candidate pool.\\
	max\_cands & Maximum number of candidate chunks considered by the selection algorithm.\\
	$k$ / $n$ & Number of chunks to select / total candidates.\\
	$c_r$ & Compression ratio (Wang et al.): fraction of tokens selected over total.\\
	\bottomrule
\end{tabular}

\end{document}
